\begin{itemize}
\item \textbf{Synthetic, easy, high ceiling.} Four relations are derived exactly by two-hop rules from facts that are present with high probability in the observed graph, and the rules are redundant. The path model is near the ceiling, so the study cannot rank path models, say anything about longer rules, noisy or non-rule facts, or realistic benchmarks, and the depth ablation cannot show an advantage of deeper models.
\item \textbf{Same-distribution new graph.} The ``fresh'' graph is generated by the same process with the same relation vocabulary; the inductive setting here is new entities only, not new relations and not a distribution shift in graph structure or size.
\item \textbf{Scale and statistics.} \rhval{const/n_seeds} seeds, graphs of \rhval{const/people} people, a few hundred test triples per seed; paired tests with five seeds detect only large effects, and we never interpret non-significance as equivalence (for example the third layer, or Path-MP versus fold-in TransE at 25\% observed).
\item \textbf{Reimplemented, reduced baselines.} TransE is our own implementation with a custom, modest negative-sampling and normalisation scheme; other tuned embedding models (RotatE, ComplEx) can represent symmetric and compositional patterns and would probably do better than TransE, so the gap between embeddings and Path-MP is an upper bound for the embedding family. We did not run GraIL, NBFNet, DRUM or ULTRA; Path-MP is a reduced NBFNet-style model. Tuning effort was larger for TransE than for Path-MP.
\item \textbf{Fold-in uses the observed new graph.} The fold-in baseline needs gradient steps on the new graph; we did not tune its number of steps separately, and we did not compare against a cheaper inductive embedding scheme.
\item \textbf{Oracle is not a ceiling.} The rule oracle applies each rule once; the learned model exceeds it, so no result here should be read as ``fraction of the achievable''.
\item \textbf{Mechanisms are hypotheses.} Explanations for the one-layer failure, the no-edge-removal failure and TransE's weakness on symmetric relations are plausible readings of the numbers, not isolated by dedicated runs. The advisory verdict tool of the pipeline compares the method with the rule oracle only and is not informative here.
\item \textbf{Split design.} Units are held out at random; triples that cannot be inferred from the observed graph exist and bound the achievable MRR (the parent relation), but we did not count them.
\end{itemize}
